\section{Arquitectura propuesta}
\label{sec:arq}
\begin{respuesta}
Dos planos y una capa transversal. El \textbf{plano de control} decide \emph{qué} puede existir
(registro, catálogos, constructor, compuerta de evaluación, políticas). El \textbf{plano de
ejecución} decide \emph{cómo} corre (API de agentes, runtime, gateway de modelos, servidores
de herramientas y conocimiento). La \textbf{observabilidad} registra todo y devuelve señales
al plano de control.
\end{respuesta}

La arquitectura está pensada para que la experiencia sea lo más fluida posible para las dos
personas que la usan: quien crea un agente (lo arma conversando, lo prueba y lo publica sin
abrir un ticket) y quien conversa con él (respuestas que fluyen, con datos reales y sin
fricción). Parte importante del proceso fue \textbf{condensar la arquitectura} para minimizar
costos, evitar piezas móviles que puedan fallar y no dejar complejidad de legado a los
desarrolladores que vengan después: un solo runtime, un solo gateway, una sola base de datos
y una sola API para todos los agentes.

Todas las decisiones que siguen salen de seis principios, y son ellos los que deberían
sobrevivir si mañana cambian las herramientas concretas:
\begin{enumerate}[label=\textbf{P\arabic*.}, leftmargin=2.2em]
  \item El agente es configuración declarativa, no código.
  \item Existe un solo camino a producción, y es el más fácil.
  \item Los agentes piden capacidades, no modelos.
  \item Lo que se recupera (documentos, herramientas) es dato, nunca instrucción.
  \item Nada se publica ni cambia sin evaluación.
  \item La solución más simple que funcione \citep{anthropic2024effective}.
\end{enumerate}

% --------------------------------------------------------------
\subsection{Contexto}
\begin{figure}[!htbp]
  \centering
  \begin{tikzpicture}[x=1cm,y=1cm, act/.style={neut, rounded corners=8pt, text width=2.6cm, minimum height=0.95cm}, ext/.style={dat, text width=2.6cm, minimum height=0.95cm}]
\node[ctrl, fill=docink, text=white, draw=docink, text width=4.6cm, minimum height=2.1cm, font=\sffamily\small] (P) at (0,0)
  {\textbf{Plataforma de agentes}\\[2pt]\scriptsize registro · catálogos · constructor\\evaluación · runtime · gateway · trazas};
% actores (izquierda)
\node[act] (a1) at (-6.2,1.6)  {\textbf{Creadores}\\servicio, cobranzas, ventas};
\node[act] (a2) at (-6.2,0)    {\textbf{Equipo de IA}\\cura catálogos y plantillas};
\node[act] (a3) at (-6.2,-1.6) {\textbf{Cumplimiento}\\aprueba riesgo alto};
\node[act] (a4) at (0,2.65) {\textbf{Clientes}\\usan agentes en la app};
% sistemas (derecha)
\node[ext] (s1) at (6.2,1.6)  {\textbf{Proveedores LLM}\\retención cero};
\node[ext] (s2) at (6.2,0)    {\textbf{Sistemas de negocio}\\facturación, medición, CRM};
\node[ext] (s3) at (6.2,-1.6) {\textbf{Documentos}\\tarifas, procedimientos, FAQ};
\node[ext] (s4) at (0,-2.65)  {\textbf{SSO corporativo}\\identidad y roles};
\draw[flecha] (a1) -- node[etq, above, sloped]{crean agentes} (P);
\draw[flecha] (a2) -- node[etq, above]{publica catálogos} (P);
\draw[flecha] (a3) -- node[etq, below, sloped]{aprueba} (P);
\draw[flecha] (a4) -- node[etq, right]{conversan} (P);
\draw[flecha] (P) -- node[etq, above, sloped]{inferencia} (s1);
\draw[flecha] (P) -- node[etq, above]{MCP} (s2);
\draw[flecha] (P) -- node[etq, below, sloped]{ingesta} (s3);
\draw[flecha] (s4) -- node[etq, right]{autentica} (P);
\end{tikzpicture}

  \caption{Vista de contexto. A la izquierda, las personas que crean, curan y aprueban; a la
  derecha, los sistemas externos. La plataforma es el único punto de contacto entre ambos.}
  \label{fig:contexto}
\end{figure}

Cuatro tipos de personas usan la plataforma con papeles distintos: los creadores de cada
área, el equipo de IA que cura catálogos y plantillas, cumplimiento que aprueba lo de riesgo
alto, y los clientes que conversan con los agentes publicados (\cref{fig:contexto}).

% --------------------------------------------------------------
\subsection{Componentes y stack}
La \cref{fig:arquitectura} muestra los componentes y la \cref{tab:stack} el stack con que los
implemento. Estas son las decisiones que los definen.

\begin{figure}[!htbp]
  \centering
  \begin{tikzpicture}[x=1cm,y=1cm]
% ---------- Consumidores ----------
\node[neut, minimum width=3.0cm, minimum height=0.75cm] (emp) at (5.0,8.15) {\textbf{Empleados}\\constructor y consola};
\node[neut, minimum width=2.4cm, minimum height=0.75cm] (can) at (12.05,8.15) {\textbf{Canales}\\chat, Slack, app de clientes};
\node[neut, minimum width=2.4cm, minimum height=0.75cm] (aut) at (14.95,8.15) {\textbf{Automatización}\\flujos, sistemas};

% ---------- Plano de control ----------
\node[ctrl, minimum width=5.6cm, minimum height=0.62cm] (cons) at (3.3,6.55) {\textbf{Constructor + Arquitecto}\\entrevista guiada, compilación};
\node[ctrl, minimum width=5.6cm, minimum height=0.62cm] (reg) at (3.3,5.6) {\textbf{Registro de agentes}\\manifiestos, versiones, estados, dueños};
\node[ctrl, text width=2.45cm, minimum height=0.75cm] (cmod) at (1.85,4.5) {\textbf{Modelos}\\alias y capacidades};
\node[ctrl, text width=2.45cm, minimum height=0.75cm] (corq) at (4.75,4.5) {\textbf{Orquestación}\\herramientas y conocimiento};
\node[ctrl, minimum width=5.6cm, minimum height=0.55cm] (gate) at (3.3,3.45) {\textbf{Compuerta de evaluación}\\suites propias y comunes, umbral por riesgo};
\node[ctrl, minimum width=5.6cm, minimum height=0.45cm] (pol) at (3.3,2.6) {Identidad (SSO) · políticas de datos · presupuestos};
\begin{scope}[on background layer]
  \node[grupo, draw=ctrlbd, fill=ctrlbg!35, fit=(cons)(reg)(cmod)(corq)(gate)(pol), inner ysep=9pt] (CP) {};
\end{scope}
\node[titulo, text=ctrlbd, anchor=south west] at (CP.north west) {PLANO DE CONTROL};

% ---------- Plano de ejecución ----------
\node[exec, minimum width=7.6cm, minimum height=0.62cm] (api) at (12.05,6.55) {\textbf{API de agentes}\\compatible con OpenAI · auth · streaming};
\node[exec, minimum width=7.6cm, minimum height=0.8cm] (rt) at (12.05,5.4) {\textbf{Runtime} (grafo construido desde el manifiesto)\\guard de entrada $\rightarrow$ agente $\leftrightarrow$ herramientas $\rightarrow$ guard de salida (por bloques)};
\node[exec, minimum width=2.4cm, minimum height=0.95cm] (gw) at (9.45,3.75) {\textbf{Gateway de}\\\textbf{modelos}\\alias · llaves · costo\\enmascara PII};
\node[exec, minimum width=2.4cm, minimum height=0.95cm] (mcp) at (12.05,3.75) {\textbf{Servidores MCP}\\lectura / escritura\\permisos por agente};
\node[exec, minimum width=2.4cm, minimum height=0.95cm] (kn) at (14.65,3.75) {\textbf{Conocimiento}\\búsqueda híbrida\\filtros de vigencia};
\begin{scope}[on background layer]
  \node[grupo, draw=execbd, fill=execbg!35, fit=(api)(rt)(gw)(mcp)(kn), inner ysep=9pt] (EP) {};
\end{scope}
\node[titulo, text=execbd, anchor=south west] at (EP.north west) {PLANO DE EJECUCIÓN};

% ---------- Externos ----------
\node[dat, minimum width=2.4cm, minimum height=0.75cm] (prov) at (9.45,1.75) {\textbf{Proveedores LLM}\\API cerrada · abiertos};
\node[dat, minimum width=4.4cm, minimum height=0.75cm] (sis) at (13.35,1.75) {\textbf{Sistemas de negocio}\\facturación · medidores · CRM · soporte};

% ---------- Observabilidad ----------
\node[obs, minimum width=16.0cm, minimum height=0.6cm] (obsv) at (8.0,0.35) {\textbf{OBSERVABILIDAD} · trazas OpenTelemetry por agente y versión · costo por agente y área · evaluación en línea · alertas e incidentes};

% ---------- Flechas ----------
\draw[flecha] (emp) -- (emp |- cons.north);
\draw[flecha] (can) -- (can |- api.north);
\draw[flecha] (aut) -- (aut |- api.north);
\draw[flecha] (api) -- (rt);
\draw[flecha] (rt.south -| gw) -- (gw);
\draw[flecha] (rt.south -| mcp) -- (mcp);
\draw[flecha] (rt.south -| kn) -- (kn);
\draw[flecha] (gw) -- (prov);
\draw[flecha] (mcp.south) -- (mcp.south |- sis.north);
\draw[flecha] (kn.south) -- (kn.south |- sis.north);
\draw[flechad, draw=ctrlbd, line width=0.9pt] (reg.east) -- node[etq, below, text=ctrlbd, yshift=-1pt]{manifiesto} (rt.west);
\draw[flechad, draw=ctrlbd, line width=0.9pt] (CP.east |- gw.west) -- node[etq, above, text=ctrlbd, pos=0.45]{alias, permisos} (gw.west);
\draw[flecha, draw=obsbd] ($(EP.south west)+(0.45,0)$) -- node[etq, text=obsbd]{trazas} ($(EP.south west)+(0.45,0)$ |- obsv.north);
\draw[flecha, draw=obsbd] (obsv.north -| gate) ++(0,0) -- (pol.south) node[etq, midway, right, text=obsbd, xshift=2pt]{señales: calidad, costo, incidentes};
\end{tikzpicture}

  \caption{Componentes. Azul: plano de control. Verde: plano de ejecución. Ámbar:
  observabilidad. Punteado: configuración; continuo: ejecución.}
  \label{fig:arquitectura}
\end{figure}

\begin{table}[!htbp]
\centering\small
\caption{Stack propuesto. Todo es de código abierto o estándar, y todo corre dentro de la
infraestructura de la empresa salvo los proveedores de modelos.}
\label{tab:stack}
\renewcommand{\arraystretch}{1.1}
\begin{tabularx}{\textwidth}{@{}L{3.6cm}L{4.4cm}Y@{}}
\toprule
\textbf{Pieza} & \textbf{Tecnología} & \textbf{Rol}\\
\midrule
Runtime de agentes & \textbf{LangGraph} (Python) & Convierte cada manifiesto en un grafo ejecutable\\
API y plano de control & FastAPI + Pydantic & API compatible con OpenAI, registro, validación del manifiesto\\
Gateway de modelos & LiteLLM & Alias, llaves, costo, reintentos y respaldo\\
Herramientas & Servidores MCP & Conectores a facturación, medición, CRM y soporte\\
Datos & PostgreSQL + pgvector & Registro, vectores, estado de conversaciones\\
Observabilidad & Langfuse + OpenTelemetry & Trazas, costos y evaluación en línea\\
Consola y constructor & Aplicación web & Inventarios, constructor conversacional, aprobaciones\\
\bottomrule
\end{tabularx}
\end{table}

\decision{Arquitectura condensada: un runtime, un gateway, una base de datos, una API.}
{Cada pieza adicional es algo más que puede fallar, que hay que pagar y que alguien tendrá que
mantener. Con estas cuatro se cubren todos los agentes de la empresa; lo demás se agrega solo cuando
un caso real lo pida.}
{Microservicios por capacidad, colas y una base de datos por función desde el primer día.}

\decision{LangGraph como framework de agentes.}
{Modela cada agente como un grafo de estados explícito, que es justo lo que produce un
manifiesto: se puede leer, probar y auditar paso a paso. Guarda el estado en PostgreSQL, así
que una conversación se reanuda después de una aprobación humana (necesario para L4), y trae
streaming nativo. Además, ya lo he llevado a producción.}
{Marcos multiagente conversacionales (el flujo queda implícito en la charla entre agentes y es
difícil de auditar) y un orquestador propio (otro framework que mantener).}

\decision{Separar plano de control y plano de ejecución.}
{Lo que cambia poco y requiere aprobación (qué agentes existen, con qué permisos) no debe
mezclarse con lo que corre miles de veces al día. Se escalan, despliegan y auditan por
separado.}
{Un servicio monolítico por agente, que es exactamente la situación actual.}

\decision{El gateway es el único componente con llaves de proveedores.}
{Así cada llamada a un modelo queda atribuida a un agente y un área, con costo y política de
datos aplicados en un solo lugar. ``¿Cuánto cuesta cada agente?'' se vuelve una consulta.}
{Llaves por equipo o por agente: es la causa de que hoy nadie sepa cuánto se gasta.}

\decision{Herramientas expuestas por MCP \citep{mcp2025spec}.}
{Cada sistema de la empresa se conecta una vez y queda disponible para todo agente autorizado, con
operaciones de lectura y escritura separadas y permisos declarados en el manifiesto.}
{Funciones embebidas en cada agente: se reimplementan, se revisan N veces y nadie sabe quién
puede escribir en qué.}

\decision{API de agentes compatible con OpenAI.}
{Cada agente publicado aparece como un ``modelo'': cualquier canal (chat interno, Slack, app
de clientes) se conecta sin integración a la medida. Ya lo he operado en producción.}
{Una API propia por canal.}

% --------------------------------------------------------------
\subsection{Una consulta de punta a punta}
\begin{figure}[!htbp]
  \centering
  \begin{tikzpicture}[x=1cm,y=1cm]
\def\yb{-7.55}
\node[obs, text width=15.4cm, align=left, anchor=west, minimum height=0.6cm] at (-0.95,1.05)
  {\textbf{Proceso intervenido: consulta de factura en servicio al cliente.} Hoy un asesor revisa
  facturación, medición y la normativa tarifaria para responder. Con la plataforma, el asistente
  responde con los datos del cliente y escala a un asesor lo que no le corresponde.};
\foreach \x/\n/\t/\s in {0/cli/Cliente/neut, 2.6/api/{API + Runtime}/exec, 5.1/grd/Guards/risk,
                         7.6/gw/Gateway/exec, 10.1/llm/{Modelo (alias)}/dat, 12.6/mcp/{MCP: facturación}/dat,
                         15.0/tr/Trazas/obs}{
  \node[\s, minimum width=2.0cm, minimum height=0.55cm, font=\sffamily\scriptsize\bfseries] (\n) at (\x,0) {\t};
  \draw[densely dotted, draw=neutbd!60, line width=0.6pt] (\n.south) -- (\x,\yb);
}
\newcommand{\msg}[5][flecha]{%
  \draw[#1] (#2,#4) -- node[etq, above, yshift=-0.5pt]{#5} (#3,#4);}
\msg{0}{2.6}{-0.75}{pregunta + sesión}
\msg{2.6}{5.1}{-1.3}{entrada: inyección, alcance}
\msg[flechad]{5.1}{2.6}{-1.8}{ok}
\msg{2.6}{7.6}{-2.35}{alias \texttt{estandar}}
\msg{7.6}{10.1}{-2.9}{modelo resuelto + política}
\msg[flechad]{10.1}{7.6}{-3.45}{pide herramienta}
\msg[flechad]{7.6}{2.6}{-3.95}{\texttt{consultar\_factura}}
\msg{2.6}{12.6}{-4.5}{cliente\_id \emph{de la sesión}, no del texto}
\msg[flechad]{12.6}{2.6}{-5.05}{datos marcados como no confiables}
\msg{2.6}{10.1}{-5.6}{segunda llamada con el resultado (vía gateway)}
\msg[flechad]{10.1}{2.6}{-6.15}{respuesta por bloques}
\msg{2.6}{5.1}{-6.65}{salida: valida cada bloque}
\msg[flechad]{2.6}{0}{-7.2}{respuesta}
% trazas
\draw[flecha, draw=obsbd, densely dashed] (7.6,-4.2) -- node[etq, above, text=obsbd, pos=0.75]{span: tokens, costo} (15.0,-4.2);
\draw[flecha, draw=obsbd, densely dashed] (2.6,-7.0) -- node[etq, above, text=obsbd, pos=0.8]{traza enmascarada} (15.0,-7.0);
% marcas de pasos
\foreach \y/\k in {-0.75/1,-1.3/2,-2.35/3,-4.5/4,-5.6/5,-6.65/6,-7.2/7}{
  \node[circle, fill=docacento, text=white, inner sep=0.8pt, font=\sffamily\tiny\bfseries] at (-0.55,\y) {\k};
}
\end{tikzpicture}

  \caption{Una consulta de factura de punta a punta. Cada flecha deja un registro en la traza
  con agente, versión, modelo, tokens, costo y latencia.}
  \label{fig:secuencia}
\end{figure}

La \cref{fig:secuencia} sigue el proceso donde la plataforma entra primero: las consultas
de factura que hoy atiende un asesor. Dos detalles sostienen la seguridad sin sacrificar la
experiencia. El identificador del cliente lo pone el runtime desde la sesión (paso 4), así
que el agente no puede consultar a otra empresa aunque se lo pidan. Y el guard de salida
valida la respuesta por bloques (paso 6): el cliente la ve aparecer mientras se escribe, pero
ningún bloque sale sin haber sido revisado.

% --------------------------------------------------------------
\subsection{Despliegue}
\begin{figure}[!htbp]
  \centering
  \begin{tikzpicture}[x=1cm,y=1cm, svc/.style={exec, text width=2.2cm, minimum height=1.0cm}, st/.style={dat, minimum height=1.0cm, rounded corners=1pt}]
% fila principal
\node[svc] (ui)  at (0,2.6)   {\textbf{Consola y constructor}\\web};
\node[svc] (api) at (3.0,2.6) {\textbf{API de agentes}\\FastAPI, réplicas};
\node[exec, text width=2.5cm, minimum height=1.0cm] (rt) at (6.15,2.6) {\textbf{Runtime}\\LangGraph, workers};
\node[svc] (gw)  at (9.3,2.6) {\textbf{Gateway}\\LiteLLM};
% arriba
\node[dat, text width=2.2cm, minimum height=0.85cm] (sec) at (9.3,4.25) {\textbf{Secretos}\\llaves solo del gateway};
% segunda fila
\node[svc] (ing) at (0,0.9)   {\textbf{Ingesta}\\conversión, troceo};
\node[svc] (ev)  at (3.0,0.9) {\textbf{Evaluación}\\jobs programados};
\node[exec, text width=2.25cm, minimum height=1.0cm] (mcp) at (6.5,0.9) {\textbf{Servidores MCP}\\uno por sistema};
\node[st, text width=2.2cm] (lf) at (9.3,0.9) {\textbf{Langfuse}\\trazas y costos};
% datos
\node[st, text width=3.0cm] (pg) at (1.5,-0.85) {\textbf{PostgreSQL}\\registro, pgvector, estado};
\begin{scope}[on background layer]
\node[grupo, draw=execbd, fill=execbg!25, fit=(ui)(gw)(sec)(pg)(lf)(mcp), inner sep=9pt] (V) {};
\end{scope}
\node[titulo, text=execbd, anchor=south west] at (V.north west) {NUBE DE LA EMPRESA · red privada};
% externos
\node[dat, text width=2.3cm, minimum height=1.0cm] (prov) at (13.3,2.6) {\textbf{Proveedores LLM}\\retención cero};
\node[dat, text width=2.3cm, minimum height=1.0cm] (sis) at (13.3,-0.85) {\textbf{Sistemas de negocio}\\facturación, medición, CRM};
\node[neut, text width=2.3cm, minimum height=0.85cm] (sso) at (13.3,4.25) {\textbf{SSO}\\identidad};
% flujo principal
\draw[flecha] (ui) -- (api); \draw[flecha] (api) -- (rt); \draw[flecha] (rt) -- (gw);
\draw[flecha, line width=1.1pt, draw=docacento] (gw) -- node[etq, above, text=docacento]{única salida} (prov);
\draw[flecha] (sec) -- (gw);
\draw[flecha] (ev) -- (api);
\draw[flecha] ($(rt.south)+(0.3,0)$) -- (mcp.north -| {$(rt.south)+(0.3,0)$});
\draw[flecha, draw=obsbd] (gw) -- node[etq, right, text=obsbd]{trazas} (lf);
% datos
\draw[flechad] (ing.south -| {$(pg.north west)+(0.35,0)$}) -- ($(pg.north west)+(0.35,0)$);
\draw[flechad] ($(rt.south west)+(0.12,0)$) |- node[etq, pos=0.75, below]{estado} (pg.east);
% externos
\draw[flecha] (mcp.south) |- (sis.west);
\draw[flecha] (sso.north) -- ++(0,0.8) -| node[etq, pos=0.25, above]{identidad} (api.north);
\end{tikzpicture}

  \caption{Despliegue en la nube de la empresa. El gateway es la única salida hacia proveedores de
  modelos; las trazas y los índices quedan dentro de la red privada.}
  \label{fig:despliegue}
\end{figure}

\decision{Todo en contenedores dentro de la red de la empresa, con PostgreSQL como almacén único
(registro, vectores con pgvector y estado de conversaciones).}
{El volumen de la empresa no justifica una base vectorial dedicada; un solo motor simplifica
respaldos, permisos y operación.}
{Base vectorial administrada aparte; plataforma de agentes de un proveedor de nube (ata
modelos al proveedor y deja los inventarios fuera de la empresa).}

\decision{Observabilidad autohospedada (Langfuse) con convenciones GenAI de OpenTelemetry
\citep{otel2025genai}.}
{Las trazas contienen conversaciones con clientes; aun enmascaradas, conviene que no salgan
de la infraestructura de la empresa \citep{ley1581}.}
{Servicio SaaS de trazas.}

% --------------------------------------------------------------
\subsection{Seguridad}
\begin{figure}[!htbp]
  \centering
  \begin{tikzpicture}[x=1cm,y=1cm]
\foreach \i/\c/\capa/\que/\riesgo in {
  0/ctrl/{1. Identidad}/{SSO; el sujeto (cliente) sale de la sesión, nunca del texto}/{consultar datos de otro cliente},
  1/risk/{2. Guard de entrada}/{clasificador de inyección y de alcance; límites de longitud}/{inyección directa},
  2/ctrl/{3. Permisos y autonomía}/{herramientas declaradas; escritura con confirmación; nivel L1--L4}/{agencia excesiva},
  3/dat/{4. Lo recuperado es dato}/{resultados de herramientas y RAG marcados como no confiables}/{inyección indirecta},
  4/risk/{5. Guard de salida por bloques}/{cifras verificadas contra herramientas; PII ajena; reglas internas}/{fuga de datos o del prompt},
  5/obs/{6. Trazas y presupuestos}/{PII enmascarada antes de persistir; tope por conversación y mes}/{exposición en logs; costo sin límite}}{
  \node[\c, text width=2.9cm, minimum height=0.62cm, anchor=west, font=\sffamily\scriptsize\bfseries] (k\i) at (0,-0.78*\i) {\capa};
  \node[plain, text width=7.6cm, anchor=west, align=left] at (3.35,-0.78*\i) {\que};
  \node[plain, text width=3.6cm, anchor=west, align=left, text=riskbd] at (11.3,-0.78*\i) {\riesgo};
}
\node[titulo, anchor=west] at (0,0.62) {CAPA};
\node[titulo, anchor=west] at (3.35,0.62) {QUÉ HACE};
\node[titulo, anchor=west, text=riskbd] at (11.3,0.62) {QUÉ DETIENE (OWASP LLM)};
\draw[-{Stealth[length=6pt]}, line width=1.2pt, draw=docacento] (-0.35,0.35) -- (-0.35,-4.25) node[midway, left, font=\sffamily\tiny, text=docacento, align=center]{orden\\en cada\\petición};
\end{tikzpicture}

  \caption{Seis capas que atraviesa cada petición y el riesgo del OWASP Top 10 para LLM
  \citep{owasp2025llm} que detiene cada una.}
  \label{fig:seguridad}
\end{figure}

\decision{Los controles son capas de la plataforma que todo agente hereda; ningún creador las
puede quitar.}
{Si cada agente implementa su propia seguridad, la seguridad depende de quién lo construyó.
La inyección indirecta, por ejemplo, llega dentro de documentos y respuestas de herramientas
\citep{greshake2023indirect} y solo se contiene tratando todo lo recuperado como dato.}
{Filtros por expresiones regulares por agente: los usé y fallaron en ambos sentidos.}

% --------------------------------------------------------------
\subsection{Clasificación de datos}
\begin{figure}[!htbp]
  \centering
  \begin{tikzpicture}[x=1cm,y=1cm, cl/.style={base, text width=3.3cm, minimum height=1.05cm, font=\sffamily\scriptsize}]
\foreach \x/\c/\n in {0/execbg/Pública, 3.75/ctrlbg/Interna, 7.5/obsbg/Confidencial, 11.25/riskbg/Restringida}{
  \node[cl, fill=\c, draw=neutbd, minimum height=0.75cm, font=\sffamily\small\bfseries] at (\x,0) {\n};}
\foreach \x/\t in {0/{tarifas publicadas, FAQ}, 3.75/{procedimientos, guías internas}, 7.5/{facturas, consumo, contacto de clientes}, 11.25/{datos bancarios, credenciales, cédulas}}{
  \node[cl, fill=white, draw=neutbd!40] at (\x,-1.15) {\t};}
\foreach \x/\t in {0/{cualquier tier del catálogo}, 3.75/{proveedores con acuerdo empresarial}, 7.5/{solo retención cero + acuerdo de datos}, 11.25/{\textbf{nunca} al modelo: se enmascara o va por herramienta}}{
  \node[cl, fill=white, draw=neutbd!40] at (\x,-2.35) {\t};}
\foreach \x/\t in {0/{completas}, 3.75/{completas}, 7.5/{PII enmascarada}, 11.25/{solo metadatos}}{
  \node[cl, fill=white, draw=neutbd!40] at (\x,-3.55) {\t};}
\foreach \y/\t in {-1.15/Ejemplos, -2.35/Modelos, -3.55/Trazas}{
  \node[titulo, anchor=east] at (-1.9,\y) {\t};}
\draw[-{Stealth[length=6pt]}, line width=1.1pt, draw=riskbd] (-1.75,0.75) -- (13.0,0.75);
\node[font=\sffamily\tiny, text=riskbd, anchor=south east] at (13.0,0.8) {más restricción};
\end{tikzpicture}

  \caption{Cuatro niveles de clasificación. El manifiesto declara el nivel del agente y el
  gateway restringe los modelos y el contenido de las trazas en consecuencia.}
  \label{fig:datos}
\end{figure}

\decision{La clasificación del dato se decide antes que el modelo.}
{Es lo único que no se negocia por calidad o costo: restringe el conjunto de modelos elegibles
y lo que se guarda en trazas.}
{Dejar que cada área escoja proveedor según preferencia.}

% --------------------------------------------------------------
\subsection{Inventarios}
Los inventarios de agentes, modelos y orquestación no son hojas de cálculo que alguien mantiene: son vistas del
plano de control que se llenan solas, porque sin registro no hay camino a producción
(\cref{tab:inventarios}).

\begin{table}[!htbp]
\centering\small
\caption{Inventarios y la pregunta que responde cada uno.}
\label{tab:inventarios}
\renewcommand{\arraystretch}{1.1}
\begin{tabularx}{\textwidth}{@{}L{2.2cm}YL{3.9cm}@{}}
\toprule
\textbf{Inventario} & \textbf{Qué registra} & \textbf{Pregunta}\\
\midrule
Agentes & dueño, área, versión, alias, herramientas, colecciones, clasificación, autonomía, estado, costo & ¿Cuántos hay y qué tocan?\\
Modelos & alias, modelo y versión, proveedor, capacidades verificadas, retención, precio & ¿Con qué corren?\\
Orquestación & herramientas (lectura o escritura, aprobación), colecciones (dueño, vigencia), flujos & ¿Qué puede hacer cada uno?\\
\bottomrule
\end{tabularx}
\end{table}

% --------------------------------------------------------------
\subsection{Modelos}
\label{sec:modelos}
Escojo el modelo por dato y por tarea, en ese orden (\cref{fig:modelos}). Y el día que sale
un modelo mejor no se cambia en ningún agente: se cambia detrás de un alias, después de
pasar por la \cref{fig:cambiomodelo}.
\begin{figure}[!htbp]
  \centering
  \begin{tikzpicture}[x=1cm,y=1cm, dd/.style={dec, font=\sffamily\tiny, inner sep=0pt, minimum width=3.1cm, minimum height=1.45cm}]
\node[neut, minimum height=0.6cm] (ini) at (-0.9,0) {Tarea\\nueva};
\node[dd] (d1) at (1.9,0) {¿Ve datos\\restringidos?};
\node[dd] (d2) at (5.5,0) {¿Ve datos de\\clientes (confidencial)?};
\node[dd] (d3) at (9.1,0) {¿Tarea estrecha y\\de alto volumen?};
\node[dd] (d4) at (12.7,0) {¿Planifica varios\\pasos o razona?};
\draw[flecha] (ini) -- (d1);
\draw[flecha] (d1) -- node[etq, above]{no} (d2);
\draw[flecha] (d2) -- node[etq, above, pos=0.3]{no} (d3);
\draw[flecha] (d3) -- node[etq, above]{no} (d4);
\node[risk, text width=2.6cm] (o1) at (1.9,-2.25) {\textbf{No va al modelo}\\se enmascara o se resuelve con herramienta};
\node[ctrl, text width=2.6cm] (o2) at (5.5,-2.25) {\textbf{Restricción}: solo proveedores con retención cero y acuerdo de datos};
\node[exec, text width=2.4cm] (o3) at (9.1,-2.25) {\textbf{Tier} \texttt{rapido}\\clasificar, extraer, resumir};
\node[exec, text width=2.75cm] (o4) at (12.55,-2.25) {\textbf{Tier} \texttt{razonamiento}\\planificación, análisis};
\node[exec, text width=2.0cm] (o5) at (15.35,-2.25) {\textbf{Tier} \texttt{estandar}\\conversación con herramientas};
\draw[flecha] (d1) -- node[etq, right]{sí} (o1);
\draw[flecha] (d2) -- node[etq, right]{sí} (o2);
\draw[flecha] (d3) -- node[etq, right]{sí} (o3);
\draw[flecha] (d4) -- node[etq, right]{sí} (o4);
\draw[flecha] (d4.east) -| node[etq, above, pos=0.25]{no} (o5.north);
\coordinate (jn) at ($(d2.east)!0.5!(d3.west)$);
\draw[dashed, line width=0.7pt, draw=neutbd] (o2.east) -| (jn);
\fill[neutbd] (jn) circle (1.6pt);
\node[plain, text=docgris, font=\sffamily\tiny\itshape, anchor=north] at (o2.south) {la restricción se suma; la decisión continúa};
\end{tikzpicture}

  \caption{Cómo se asigna un tier: primero el dato, después la tarea.}
  \label{fig:modelos}
\end{figure}

\decision{Tres alias (\texttt{rapido}, \texttt{estandar}, \texttt{razonamiento}) en lugar de
nombres de modelo; tier fijo por agente al inicio.}
{Las tareas estrechas no necesitan el modelo más caro \citep{chen2023frugalgpt,ong2025routellm}.
El enrutamiento dinámico ahorra más, pero agrega una pieza que evaluar: lo dejo para la
fase 3.}
{Un único modelo para todo (caro) o enrutamiento dinámico desde el día uno (complejo).}

\decision{Modelos cerrados por API con retención cero por defecto; sin GPU propias.}
{Hoy dan la mejor capacidad por esfuerzo operativo. Lo local solo se justifica si el dato no
puede salir o si el volumen hace más barato operar GPU, y ninguna condición se cumple en la empresa
al inicio. Los modelos abiertos entran al catálogo cuando ganan en costo para un tier.}
{Desplegar modelos abiertos propios desde el inicio.}

\begin{figure}[!htbp]
  \centering
  \begin{tikzpicture}[x=1cm,y=1cm, st/.style={exec, text width=2.05cm, minimum height=1.25cm}]
\node[dat, text width=1.7cm, minimum height=1.25cm] (s0) at (0,0) {\textbf{Candidato}\\sale un modelo nuevo};
\node[st] (s1) at (2.75,0) {\textbf{1. Capacidades}\\herramientas, JSON, contexto, español};
\node[st] (s2) at (5.55,0) {\textbf{2. Suites offline}\\de todos los agentes del alias};
\node[st] (s3) at (8.35,0) {\textbf{3. Sombra}\\tráfico real, sin responder; se compara};
\node[st] (s4) at (11.15,0) {\textbf{4. Canario}\\5--10\,\% del tráfico con monitoreo};
\node[ctrl, text width=1.9cm, minimum height=1.25cm] (s5) at (13.95,0) {\textbf{5. Mover alias}\\anterior queda para reversión};
\foreach \a/\b in {s0/s1,s1/s2,s2/s3,s3/s4,s4/s5}{\draw[flecha] (\a) -- (\b);}
\node[risk, text width=9.4cm] (rej) at (6.95,-2.05) {\textbf{Rechazo registrado}: el resultado queda en el catálogo; ningún agente cambia};
\foreach \s in {s1,s2,s3,s4}{\draw[flechad, draw=riskbd] (\s.south) -- (\s.south |- rej.north);}
\node[etq, text=riskbd] at (4.15,-1.2) {falla};
\node[obs, text width=3.0cm] (mon) at (13.95,-2.05) {Monitoreo en producción; reversión en un paso};
\draw[flecha] (s5) -- (mon);
\end{tikzpicture}

  \caption{Qué pasa el día que sale un modelo mejor: cada etapa tiene criterio de salida y
  reversión inmediata.}
  \label{fig:cambiomodelo}
\end{figure}

\decision{El catálogo registra capacidades \emph{verificadas} (herramientas, JSON válido,
contexto efectivo, español) y fija versiones concretas.}
{Un nombre de modelo no garantiza capacidades: lo viví cuando un cambio de modelo rompió la salida
estructurada de una función que nadie volvió a probar. Además, un mismo modelo comercial cambia de comportamiento entre versiones
\citep{chen2024drift}.}
{Confiar en lo que anuncia el proveedor y usar alias ``latest''.}

% --------------------------------------------------------------
\subsection{Conocimiento, RAG y fine-tuning}
\label{sec:conocimiento}
RAG \citep{lewis2020rag} sirve cuando el conocimiento es extenso, cambia y hay que citarlo.
Pero el error más común que veo es usarlo para todo. Antes de llegar a RAG descarto cuatro
casos (\cref{fig:conocimiento}): datos de un cliente, cálculos, reglas que siempre aplican y
conocimiento pequeño y estable. Para los documentos que sí van a RAG, la ingesta y la consulta
siguen la \cref{fig:ingesta}.
\begin{figure}[!htbp]
  \centering
  \begin{tikzpicture}[x=1cm,y=1cm, dd/.style={dec, font=\sffamily\tiny, inner sep=0pt, minimum width=2.75cm, minimum height=1.35cm}]
\node[neut] (q) at (-0.6,0) {Conocimiento\\que el agente\\necesita};
\node[dd] (d1) at (2.35,0) {¿Es de un cliente\\o cambia en vivo?};
\node[dd] (d2) at (5.55,0) {¿Es un cálculo?};
\node[dd] (d3) at (8.75,0) {¿Es una regla\\que siempre aplica?};
\node[dd] (d4) at (11.95,0) {¿Es pequeño\\y estable?};
\node[exec, text width=1.9cm, minimum height=1.0cm] (rag) at (15.1,0) {\textbf{RAG}\\extenso, cambia, hay que citar};
\draw[flecha] (q) -- (d1);
\foreach \a/\b in {d1/d2,d2/d3,d3/d4,d4/rag}{\draw[flecha] (\a) -- node[etq, above]{no} (\b);}
\node[dat, text width=2.5cm] (o1) at (2.35,-2.1) {\textbf{Herramienta}\\consulta al sistema fuente (factura, consumo)};
\node[dat, text width=2.5cm] (o2) at (5.55,-2.1) {\textbf{Herramienta determinista}\\el modelo no hace aritmética};
\node[risk, text width=2.5cm] (o3) at (8.75,-2.1) {\textbf{Política o guard}\\no se ``recupera'', se impone};
\node[ctrl, text width=2.5cm] (o4) at (11.95,-2.1) {\textbf{En las instrucciones}\\pocas FAQ, tono, glosario};
\foreach \d/\o in {d1/o1,d2/o2,d3/o3,d4/o4}{\draw[flecha] (\d) -- node[etq, right]{sí} (\o);}
\end{tikzpicture}

  \caption{Dónde vive cada tipo de conocimiento: se recorre de izquierda a derecha y se
  detiene en el primer ``sí''.}
  \label{fig:conocimiento}
\end{figure}

\begin{figure}[!htbp]
  \centering
  \begin{tikzpicture}[x=1cm,y=1cm, s/.style={exec, text width=2.3cm, minimum height=1.0cm}, q/.style={ctrl, text width=2.3cm, minimum height=1.0cm}]
\node[titulo, anchor=west, text=execbd] at (-1.2,1.0) {INGESTA (cuando cambia un documento)};
\node[dat, text width=2.0cm, minimum height=1.0cm] (f) at (0,0) {\textbf{Fuentes}\\PDF, Word, web, tablas};
\node[s] (i1) at (2.9,0) {\textbf{Conversión}\\texto + OCR si hace falta};
\node[s] (i2) at (5.9,0) {\textbf{Metadatos}\\dueño, vigencia, clasificación};
\node[s] (i3) at (8.9,0) {\textbf{Troceo por tipo}\\sección, par P/R, tabla};
\node[dat, text width=2.3cm, minimum height=1.0cm] (ix) at (11.9,0) {\textbf{Índices}\\BM25 + vectores (pgvector)};
\foreach \a/\b in {f/i1,i1/i2,i2/i3,i3/ix}{\draw[flecha] (\a) -- (\b);}
\node[titulo, anchor=west, text=ctrlbd] at (-1.2,-1.45) {CONSULTA (en cada pregunta)};
\node[neut, text width=2.0cm, minimum height=1.0cm] (p) at (0,-2.45) {\textbf{Pregunta}\\del agente};
\node[q] (q1) at (2.9,-2.45) {\textbf{Filtro}\\vigencia y permisos \emph{antes} de buscar};
\node[q] (q2) at (5.9,-2.45) {\textbf{Búsqueda híbrida}\\palabras clave + semántica};
\node[q] (q3) at (8.9,-2.45) {\textbf{Fusión RRF}\\y top-5};
\node[exec, text width=2.3cm, minimum height=1.0cm] (q4) at (11.9,-2.45) {\textbf{Al agente}\\como dato, con cita};
\foreach \a/\b in {p/q1,q1/q2,q2/q3,q3/q4}{\draw[flecha] (\a) -- (\b);}
\draw[flechad] (ix.south) -- ++(0,-0.55) -| (q2.north);
\end{tikzpicture}

  \caption{Pipeline de conocimiento para tarifas, procedimientos y preguntas frecuentes.}
  \label{fig:ingesta}
\end{figure}

\decision{Cifras tarifarias en tablas consultables por herramienta; explicaciones y
procedimientos en RAG con vigencia obligatoria; troceo distinto por tipo de documento.}
{Una tarifa vencida no es ``menos relevante'': es incorrecta, así que se filtra antes de
buscar. La búsqueda híbrida con fusión RRF \citep{cormack2009rrf} encuentra tanto códigos
exactos como preguntas conceptuales, y pocos fragmentos precisos rinden más que muchos
\citep{liu2024lost}.}
{RAG para todo, incluidos saldos y tarifas.}

\begin{figure}[!htbp]
  \centering
  \begin{tikzpicture}[x=1cm,y=1cm]
\def\R{2.25}
\node[titulo, text=docacento, font=\sffamily\bfseries\small] at (0,0) {Volante\\de datos};
\foreach \ang/\n/\s/\t in {90/n1/exec/{\textbf{Agente en producción}\\versión certificada},
                          18/n2/obs/{\textbf{Trazas + correcciones}\\pares ``antes y después''},
                          -54/n3/dat/{\textbf{Conjunto curado}\\etiquetado y sin PII},
                          -126/n4/ctrl/{\textbf{Suites más ricas}\\casos reales de falla},
                          162/n5/ctrl/{\textbf{Mejora}\\instrucciones, ejemplos o,\\con señal, fine-tuning}}{
  \node[\s, text width=2.75cm] (\n) at ({3.7*cos(\ang)},{2.05*sin(\ang)}) {\t};
}
\foreach \a/\b in {n1/n2,n2/n3,n3/n4,n4/n5,n5/n1}{
  \draw[flecha, line width=0.9pt, draw=docacento] (\a) to[bend left=18] (\b);
}
\node[risk, text width=3.0cm, font=\sffamily\tiny] (sen) at (-7.1,1.35) {\textbf{Señal para fine-tuning}\\techo de calidad con instrucciones \emph{y} costo relevante \emph{y} miles de pares validados};
\draw[flechad, draw=riskbd] (sen.east) -- (n5.west);
\end{tikzpicture}

  \caption{El volante de datos: las correcciones humanas mejoran primero las suites y solo
  con señal clara alimentan un fine-tuning.}
  \label{fig:volante}
\end{figure}

\decision{No hacer fine-tuning en las fases 0 a 2.}
{No hay datos validados; el conocimiento del negocio cambia cada mes y el fine-tuning inyecta
conocimiento peor que la recuperación \citep{ovadia2024ftrag,gekhman2024newknowledge}; y ata
la plataforma a un modelo base (choca con P3). Tendría lugar en una tarea estrecha y de alto
volumen, como resumir llamadas de ventas, destilando a un modelo pequeño
\citep{hu2022lora,hsieh2023distilling} cuando haya techo de calidad, costo relevante y miles
de pares validados.}
{Ajustar un modelo propio desde el inicio.}

% --------------------------------------------------------------
\subsection{Orquestación}
\label{sec:orquestacion}
\begin{figure}[!htbp]
  \centering
  \begin{tikzpicture}[x=1cm,y=1cm, pn/.style={grupo, draw=neutbd!70, minimum width=5.05cm, minimum height=4.6cm}]
% ---- Panel A
\node[pn] (A) at (0,0) {};
\node[titulo, anchor=north] at (A.north) {(a) AGENTE ÚNICO + HERRAMIENTAS};
\node[exec, minimum width=1.9cm, minimum height=0.6cm] (ag) at (0,0.75) {\textbf{Agente}\\razona y actúa};
\node[dat, font=\sffamily\tiny] (t1) at (-1.55,-0.6) {herramienta};
\node[dat, font=\sffamily\tiny] (t2) at (0,-0.6) {herramienta};
\node[dat, font=\sffamily\tiny] (t3) at (1.55,-0.6) {conocimiento};
\foreach \t in {t1,t2,t3}{\draw[flecha, {Stealth[length=4pt]}-{Stealth[length=4pt]}] (ag) -- (\t);}
\node[plain, text width=4.6cm, font=\sffamily\tiny, anchor=south] at ($(A.south)+(0,0.12)$) {\textbf{Cuándo:} un alcance de permisos, un tier, evaluable como un todo. \textbf{Por defecto.}};
% ---- Panel B
\node[pn] (B) at (5.45,0) {};
\node[titulo, anchor=north] at (B.north) {(b) FLUJO DETERMINISTA};
\node[ctrl, font=\sffamily\tiny] (b1) at (5.45,1.15) {paso 1 · código};
\node[exec, font=\sffamily\tiny] (b2) at (5.45,0.35) {paso 2 · LLM acotado};
\node[ctrl, font=\sffamily\tiny] (b3) at (5.45,-0.45) {paso 3 · validación};
\node[neut, font=\sffamily\tiny] (b4) at (5.45,-1.2) {aprobación humana};
\draw[flecha] (b1) -- (b2); \draw[flecha] (b2) -- (b3); \draw[flecha] (b3) -- (b4);
\node[plain, text width=4.6cm, font=\sffamily\tiny, anchor=south] at ($(B.south)+(0,0.12)$) {\textbf{Cuándo:} proceso regulado o irreversible. \textbf{El código decide el flujo.}};
% ---- Panel C
\node[pn] (C) at (10.9,0) {};
\node[titulo, anchor=north] at (C.north) {(c) SUPERVISOR + ESPECIALISTAS};
\node[exec, minimum width=1.8cm] (sup) at (10.9,1.05) {\textbf{Supervisor}};
\node[exec, font=\sffamily\tiny] (e1) at (9.55,-0.3) {especialista A};
\node[exec, font=\sffamily\tiny] (e2) at (12.25,-0.3) {especialista B};
\draw[flecha] (sup) -- node[etq, left, xshift=-2pt]{traspaso\\tipado} (e1);
\draw[flecha] (sup) -- (e2);
\draw[flechad] (e1) -- (e2);
\node[plain, text width=4.6cm, font=\sffamily\tiny, anchor=south] at ($(C.south)+(0,0.12)$) {\textbf{Cuándo:} permisos o tiers distintos por parte, exploración de bajo riesgo. \textbf{El LLM decide.}};
% ---- preferencia
\draw[{Stealth[length=6pt]}-{Stealth[length=6pt]}, line width=1.1pt, draw=docacento] (-2.4,-2.7) -- (13.3,-2.7);
\node[etq, text=docacento, font=\sffamily\scriptsize, anchor=north west] at (-2.4,-2.78) {más simple y auditable};
\node[etq, text=docacento, font=\sffamily\scriptsize, anchor=north east] at (13.3,-2.78) {más flexible y difícil de verificar};
\end{tikzpicture}

  \caption{Tres patrones en orden de preferencia.}
  \label{fig:orquestacion}
\end{figure}

\decision{Agente único por defecto; dividir solo por permisos, tiers o evaluación distintos.
En procesos regulados el flujo lo decide el código; los traspasos son objetos tipados, nunca
el historial completo.}
{En más de 1.600 trazas de siete marcos multiagente, las fallas vinieron sobre todo de
especificación, desalineación entre agentes y verificación, no del modelo
\citep{cemri2025mast}. El patrón razonar y actuar con herramientas \citep{yao2023react} cubre
casi todos los casos de la empresa.}
{Arquitectura multiagente desde el inicio con un supervisor LLM decidiendo todo.}
